% diagram-kind: memory
\begin{tikzpicture}[font=\small\sffamily]
\node[block,text width=27mm] (scale) at (1.5,0) {Scale bytes\\\texttt{00 38}\\binary16 $d=0.5$};
\node[block,text width=70mm] at (7.5,0) {16 packed bytes\\\texttt{F0 E1 D2 \ldots\ 1E 0F}\\two nibbles per byte};
\node[anchor=west,font=\bfseries\sffamily] at (0,-1.6) {Zoom into the first packed byte: \texttt{F0}};
\foreach \i/\bit in {0/1,1/1,2/1,3/1,4/0,5/0,6/0,7/0}
 {\node[cell,minimum width=10mm,minimum height=10mm] at (2.3+\i,-2.5) {\bit};}
\draw[BookBlue,thick] (1.8,-3.15)--(1.8,-3.45)--(5.8,-3.45)--(5.8,-3.15);
\draw[BookBlue,thick] (5.8,-3.15)--(5.8,-3.45)--(9.8,-3.45)--(9.8,-3.15);
\node[sysbox,fill=BookGreen] (hi) at (3.8,-4.7) {High nibble = 15\\$w_{16}=0.5(15-8)=3.5$};
\node[sysbox,fill=BookGold] (lo) at (8.3,-4.7) {Low nibble = 0\\$w_0=0.5(0-8)=-4$};
\draw[flow] (3.8,-3.45)--(hi);
\draw[flow] (7.8,-3.45)--(lo);
\node[sysnote,anchor=north west] at (0,-6.15) {Low nibbles produce positions 0--15. High nibbles produce positions 16--31. Do not interleave the reconstructed values.};
\end{tikzpicture}
